\chapter{Introduction and General Information} \label{intro}


This manual describes the input required by the EIRENE
code to run a Monte Carlo study for a fully 3 dimensional simulation of
linear transport (i.e., of test particles) in a prescribed background
medium. Although the geometry of the problem and the interaction between
test particle species and the background are in principle not subject to
any restrictions, the aim of code development was to provide a tool
for investigating neutral gas transport in tokamak plasmas. Consequently
the choice of preprogrammed options has been made mainly with this
application in mind. A large variety of problems in this field can be
run
without having to resort to any problem specific routines but instead by
an appropriate setting of logical and numerical input flags. \\
Any user of EIRENE should be aware that this
 code is a moving target as is this manual.\\
Therefore it is possible that there are some inconsistencies
between this description and a particular version of the code.
The user should always
first check subroutine INPUT, where most of the data
are read in using hard wired formats, as most input errors will
lead to a rapid exit in the initialization phase of a run.\\
This manual was written by the author of the code
who often may not have been
able to anticipate difficulties in understanding the use of EIRENE. He
therefore gratefully acknowledges any suggestions to make this manual
more informative and clear than it might be at the present status.\\
This code description consists of the following parts:
\begin{itemize}
\item
In the first part an introduction is given to the
general linear transport problem and its solution by Monte Carlo methods.
Most of the terminology used in later sections is introduced there.
\item
In part two a description of the formatted input file required by
EIRENE to run on a specific problem is given. It mainly consists of
explanations of the meanings of the various input flags.
\item In part three the most important problem specific routines
are explained. At present we have restricted this part to routines
for evaluation of ``user requested tallies", namely the
subroutines UPTUSR, UPCUSR and UPSUSR. Other often applied
routines such as SAMUSR (user supplied source sampling
distribution) or REFUSR (user supplied surface
interaction model) are briefly described.
\item
In part four
the package EIRCOP for
interfacing
EIRENE with other codes (e.g., the B2-EIRENE package) is
described. Here mainly the location of the storage on
the EIRENE work array RWK for plasma data and geometrical information
is given. Also the implementation of the method of (semi-implicit)
corrections (see \cref{sec1.7}) at each plasma code time-step to
the terms
transferred from EIRENE to plasma fluid transport codes is described
here.
\end{itemize}
